This section collects structural properties and local certificates used by
the two algorithms.
Monotonicity and mass bounds justify deterministic continuation in
\cref{sec:deterministic-method}. Safe active-set expansion and a local
objective certificate support the randomized method in
\cref{sec:randomized-method}. An electrical-flow interpretation is given
in \cref{app:grounded-flow-dual}.

\subsection{Nonnegative formulation and monotonicity}

On nonnegative vectors, the weighted $\ell_1$ penalty is linear. Absorb it
into the load by defining
\begin{equation}
 \vc_\rho:=\vb-\alpha\rho\vomega,
 \qquad
 \phi_\rho(\vx):=\frac12\vx^\trans\mQ\vx-\vc_\rho^\trans\vx.
 \label{eq:shifted-obstacle}
\end{equation}
Thus $F_\rho(\vx)=\phi_\rho(\vx)$ for $\vx\geq\vzero$.
Only the seed coordinate of $\vc_\rho$ can be positive, which will imply
connectivity of the optimal support.
Principal submatrices inherit the spectral bounds and off-diagonal signs
of $\mQ$, so the same Neumann-series argument gives nonnegative principal
inverses.

The reformulation and KKT conditions below follow
\citet[Section~2]{martinezrubio2023accelerated}, building on the
nonnegativity analysis of \citet[Theorem~1]{fountoulakis2019variational}.
The latter paper also discusses connected support for connected seed sets
(p.~567). We include a short proof in our point-source setting for
completeness.

\begin{lemma}[Nonnegative reformulation and connected support]
\label{thm:obstacle-reduction}
The RPPR problem has the same unique optimum and optimum value as
\begin{equation}
 \min_{\vx\geq\vzero}\phi_\rho(\vx).
 \label{eq:obstacle-problem}
\end{equation}
Its KKT system is
\begin{equation}
 \vx\geq\vzero,
 \qquad
 \vw:=\mQ\vx-\vc_\rho\geq\vzero,
 \qquad
 x_iw_i=0\quad(i\in[n]).
 \label{eq:obstacle-kkt}
\end{equation}
In the nontrivial regime \eqref{eq:nonzero-rppr-regime}, the support
$\mathcal S_\rho^*$ is connected and contains $v$.
\end{lemma}

\begin{proof}
The signs $Q_{ij}\leq0$ for $i\neq j$ and $\vb\geq\vzero$ give
$F_\rho(|\vx|)\leq F_\rho(\vx)$. Uniqueness therefore implies
$\vx_\rho^*\geq\vzero$. Equality of the objectives on the nonnegative
orthant proves the reformulation. Its optimality conditions are
\cref{eq:obstacle-kkt}, and strong convexity gives uniqueness.

At the seed,
\[
 (\vc_\rho)_v=\frac{\alpha}{\sqrt{d_v}}(1-\rho d_v)>0.
\]
If $(\vx_\rho^*)_v=0$, its slack is strictly negative, contradicting
\cref{eq:obstacle-kkt}.  Thus $v\in\mathcal S_\rho^*$.  If an active connected
component $\gC$ did not contain $v$, no support edge would join $\gC$ to the
other active components.  Stationarity on $\gC$ would give
$\mQ_{\gC\gC}(\vx_\rho^*)_\gC=(\vc_\rho)_\gC<\vzero$.  Multiplication by
$\mQ_{\gC\gC}^{-1}\geq\vzero$ contradicts positivity on $\gC$.
\end{proof}

\begin{lemma}[Order, mass, and support volume]
\label{lem:rppr-order-mass}
If $\vz\geq\vzero$ and $\mQ\vz\geq\vc_\rho$, then
$\vx_\rho^*\leq\vz$. Thus $\vx_\rho^*$ is the least nonnegative vector
satisfying this shifted-load inequality. Consequently, for $0<r\leq r'$,
\begin{equation}
 \vzero\leq\vx_{r'}^*\leq\vx_r^*\leq\vx_0^*,
 \qquad \vomega^\trans\vx_r^*+r\vol(\mathcal S_r^*)\leq1.
 \label{eq:rppr-order-mass}
\end{equation}
In particular $\vol(\mathcal S_r^*)\leq1/r$.
\end{lemma}
\begin{proof}
For the first assertion let $\gI=\{i:(\vx_\rho^*)_i>z_i\}$ and
$\ve=\vx_\rho^*-\vz$. If $\gI$ is nonempty, stationarity on
its positive optimum coordinates gives
$\mQ_{\gI\gI}\ve_{\gI}
 \leq-\mQ_{\gI\gI^c}\ve_{\gI^c}\leq\vzero$.
The last inequality uses $\ve_{\gI^c}\leq\vzero$ and Stieltjes signs.
Multiplication by the nonnegative principal inverse contradicts
$\ve_{\gI}>\vzero$. Since
$\mQ\vx_r^*\geq\vc_r\geq\vc_{r'}$ and
$\mQ\vx_0^*=\vb\geq\vc_r$, the first assertion proves the order.

The source deficit $\vb-\mQ\vx_r^*$ is nonnegative: it equals
$\alpha r\omega_i$ on the positive support, and at a zero coordinate
the off-diagonal signs give $(\mQ\vx_r^*)_i\leq0\leq b_i$.
Its weighted mass is
$\alpha(1-\vomega^\trans\vx_r^*)$, by
$\mQ\vomega=\alpha\vomega$ and $\vomega^\trans\vb=\alpha$.
Its contribution on the positive support alone is
$\alpha r\vol(\mathcal S_r^*)$, proving the mass inequality.
\end{proof}

\subsection{Safe active-set expansion}
\label{subsec:reachable-faces}

We now describe the exact active-set steps underlying the randomized
method. Assume the nonzero point-source regime in
\cref{eq:nonzero-rppr-regime}. For an active set $\gU\subseteq[n]$, define
the restricted linear-system solution, extended by zero, and its slack by
\begin{equation*}
 x_i^\gU=
 \begin{cases}
  (\mQ_{\gU\gU}^{-1}(\vc_\rho)_\gU)_i,&i\in \gU,\\
  0,&i\notin \gU,
 \end{cases}
 \qquad
 \vw^\gU:=\mQ\vx^\gU-\vc_\rho.
\end{equation*}
Set $\vx^\emptyset=\vzero$. For an arbitrary $\gU$, the restricted
solution need not be nonnegative. An active set is \emph{reachable} if it
is obtained from $\{v\}$ by repeatedly adding one or more coordinates
whose exact current slack is negative. The next theorem specializes the
restricted-minimizer geometry of
\citet[Proposition~2 and Section~3]{martinezrubio2023accelerated}.
We give a direct proof using Stieltjes inverses and Schur complements.

\begin{theorem}[Safe active-set expansion]
\label{thm:safe-batched-pivots}
For every reachable active set $\gU$,
\begin{equation*}
 \gU\subseteq\mathcal S_\rho^*,
 \qquad
 \vzero<\vx_\gU^\gU\leq(\vx_\rho^*)_\gU.
\end{equation*}
Every $j\notin \gU$ with $w_j^\gU<0$ lies in
$\mathcal S_\rho^*\setminus \gU$.  If $\gB$ is any nonempty collection of such
coordinates, then $\gU\cup \gB$ is reachable and
\begin{equation*}
 \vx_{\gU\cup \gB}^{\gU\cup \gB}>\vzero,
 \qquad
 \vx_\gU^{\gU\cup \gB}\geq\vx_\gU^\gU.
\end{equation*}
If no outside slack is negative, $\vx^\gU=\vx_\rho^*$.
\end{theorem}

\begin{proof}
The initial seed slack is negative, so the first scalar solution is
positive. Suppose the claim holds at $\gU$. Restricted stationarity and the
optimal equations give
\begin{equation*}
 \mQ_{\gU\gU}\bigl((\vx_\rho^*)_\gU-\vx_\gU^\gU\bigr)
 =-\mQ_{\gU,\mathcal S_\rho^*\setminus \gU}
   (\vx_\rho^*)_{\mathcal S_\rho^*\setminus \gU}\geq\vzero.
\end{equation*}
Inverse positivity proves the coordinatewise comparison.  If
$j\notin\mathcal S_\rho^*$, then
\begin{align*}
 w_j^\gU-w_j^*
 &=\mQ_{j\gU}\bigl(\vx_\gU^\gU-(\vx_\rho^*)_\gU\bigr)
   -\mQ_{j,\mathcal S_\rho^*\setminus \gU}
     (\vx_\rho^*)_{\mathcal S_\rho^*\setminus \gU}\geq0.
\end{align*}
Hence a negative violator cannot be a false support coordinate.

For $\gU'=\gU\cup \gB$, block elimination gives
\begin{equation*}
 \mS_\gB\vx_\gB^{\gU'}=-\vw_\gB^\gU,
 \qquad
 \mS_\gB:=\mQ_{\gB\gB}-\mQ_{\gB\gU}\mQ_{\gU\gU}^{-1}\mQ_{\gU\gB}.
\end{equation*}
Moreover,
\begin{equation*}
 \vx_\gU^{\gU'}-\vx_\gU^\gU=-\mQ_{\gU\gU}^{-1}\mQ_{\gU\gB}\vx_\gB^{\gU'}.
\end{equation*}
The Schur complement is Stieltjes and has a nonnegative inverse with
positive diagonal.  Since
$-\vw_\gB^\gU>\vzero$, the new coordinates are positive; the second display and
the off-diagonal sign make every old correction nonnegative.  This closes the
induction.  If no outside slack is negative, inside stationarity, positivity,
and outside nonnegativity satisfy the global KKT system, whose solution is
unique.
\end{proof}

Thus $\vx^\gU$ minimizes $\phi_\rho$ over nonnegative vectors supported
on a reachable $\gU$. We also call this restricted computation a
\emph{face solve}.

For any $\vx\geq\vzero$, set $\vw=\mQ\vx-\vc_\rho$ and
$y_i=x_i/\sqrt{d_i}$. At a coordinate $j\neq v$ with $x_j=0$,
\begin{equation}
 \frac{w_j}{\sqrt{d_j}}
 =\alpha\rho-\frac{1-\alpha}{2d_j}
   \sum_{i\in\mathcal N(j)\cap\supp(\vx)}y_i.
 \label{eq:boundary-slack}
\end{equation}
This expression makes support discovery local.

\begin{corollary}[Boundary-only discovery]
\label{cor:boundary-only-discovery}
At a reachable active set $\gU$, every negative outside slack lies in
$\partial\gU$. Exposing each admitted vertex's adjacency list once costs
$\vol(\gU)$ entry inspections and creates at most $\vol(\gU)$ boundary
incidences. No adjacency list outside $\gU$ is needed to evaluate the
boundary slacks.
\end{corollary}

\begin{proof}
If $j\notin \gU\cup\partial \gU$ and $j\neq v$, sparsity gives
$w_j^\gU=-(\vc_\rho)_j=\alpha\rho\sqrt{d_j}>0$.  The seed is already active.
The incidence count charges each exposed edge to its endpoint in $\gU$.
Boundary degrees can be queried without opening their adjacency lists.
Any later rescan is charged again under \cref{eq:total-work}.
\end{proof}

\subsection{Local objective certification}

We next turn slack values into a stopping test for
\cref{eq:rppr-output-target}. Let $I_{\R_+^n}$ be the indicator function,
equal to zero on $\R_+^n$ and $+\infty$ elsewhere.
For feasible $\vx\geq\vzero$, let $\vw=\mQ\vx-\vc_\rho$ and define
\begin{equation}
 g_i(\vx)=
 \begin{cases}
  w_i,&x_i>0,\\
  \min\{w_i,0\},&x_i=0.
 \end{cases}
 \label{eq:minimal-kkt-subgradient}
\end{equation}

\begin{lemma}[KKT subgradient certificate]
\label{lem:kkt-gap-certificate}
The vector $\vg(\vx)$ has the smallest Euclidean norm among the
subgradients of $\phi_\rho+I_{\R_+^n}$ at $\vx$, and
\begin{equation*}
 F_\rho(\vx)-F_\rho(\vx_\rho^*)
 \leq\frac{\norm{\vg(\vx)}_2^2}{2\alpha}.
\end{equation*}
If $\vx$ is supported on a known set $\gU$ containing $v$, the certificate
and its squared norm can be computed with $\vol(\gU)$ adjacency-entry
inspections and $\softO(\vol(\gU))$ total work under
\cref{eq:total-work}.
\end{lemma}

\begin{proof}
At $x_i>0$ the normal cone is zero.  At $x_i=0$ it is $(-\infty,0]$, so
\cref{eq:minimal-kkt-subgradient} is the closest point to zero in the
coordinate subdifferential.  Strong convexity gives, for every subgradient
$\vg$,
\[
 (\phi_\rho+I_{\R_+^n})(\vy)\geq(\phi_\rho+I_{\R_+^n})(\vx)
 +\vg^\trans(\vy-\vx)+\frac\alpha2\norm{\vy-\vx}_2^2.
\]
Taking $\vy=\vx_\rho^*$ and bounding the last two terms below by
$-\norm{\vg}_2^2/(2\alpha)$ proves the objective bound.

One scan of the rows in $\gU$ accumulates $\mQ\vx$ on
$\gU\cup\partial\gU$. Every farther coordinate has positive slack by
\cref{eq:boundary-slack}, so its certificate entry is zero.
There are at most $|\gU|+\vol(\gU)\leq2\vol(\gU)$ records.
Degree queries, diagonal and source contributions, clipping, the squared
norm, and output take $\mathcal{O}(\vol(\gU))$ elementary operations.
Accumulating entries in a balanced search tree adds a factor
$\mathcal{O}(\log(2+\vol(\gU)))$, giving the stated total work.
Degree densities permit the same computation without square roots:
the squared norm is a degree-weighted sum of squared certificate densities.
\end{proof}

In particular, $\norm{\vg(\vx)}_2^2\leq2\alpha\epsobj$ certifies the
requested RPPR accuracy using only the active set and its boundary.
